\begin{frame}[plain, noframenumbering]
	\begin{center}
		\Huge \textcolor{blue}{\textbf{Proposed Hypotheses}} \\
	\end{center}
\end{frame}

\subsection{Hypothesis 1}
\begin{frame}{Hypothesis 1 : The Moisture-Recycling Efficiency [1/2]}
	\vspace{-0.1cm}
	\begin{block}{The longevity of an MD depends on the phase-locking of the \texttt{TP-SWVL1-LHF} loop. Long-lived systems exhibit positive lags and high soil memory, whereas short-lived systems are dominated by runoff and negative lags}\end{block}

	\vspace{0.15cm}
	\footnotesize
	\textbf{Statement :} Longevity depends on the synchrony (optimal $0 \pm 6$ h, FSI $>0.4$) of the \texttt{TP-SWVL1-LHF} feedback loop. Long-lived MDs ($\ge 96$h) maintain surface energy exchange through continuous moisture recycling. Short-lived systems ($< 96$h) undergo rapid decay as land-surface friction and moisture depletion collapse the vortex core.

	\vspace{0.2cm}
	\footnotesize
	Longevity is proportional to the strength and sign of the \textbf{Land-Atmosphere Coupling}. Positive feedback (\texttt{TP}$\uparrow \rightarrow$\texttt{SWVL1}$\uparrow \rightarrow$ \texttt{LHF}$\uparrow \rightarrow$ Convection$\uparrow$) creates a \textbf{"Secondary Pseudo-Ocean"} effect. \cite{hunt2017effect}

	\vspace{0.2cm}
	\footnotesize
	\textbf{The Feedback Loop :}
	\[ \text{(\textbf{Precipitation})} \rightarrow \text{(\textbf{Infiltration})} \rightarrow \text{(\textbf{Elevated Soil Moisture})} \rightarrow \text{(\textbf{Sustained \texttt{LHF}})} \rightarrow \text{(\textbf{MD Persistence})} \]

	\vspace{0.2cm}
	\footnotesize
	\textbf{Loop Failure (Tracks 1 \& 4) :}
	\begin{itemize}
		\item \textbf{Broken Lags (\texttt{> 12h}) :} Indicates \textbf{Topographic Runoff} (Eastern Ghats) or \textbf{Saturation Excess} (Sandy Coast).
		\item \textbf{Negative Memory (\texttt{ACF $\approx -0.15$}) :} Prevents the soil from "storing" the energy recharge from precipitation.
	\end{itemize}
\end{frame}

\begin{frame}{Hypothesis 1 : The Moisture-Recycling Efficiency [2/2]}
	\footnotesize
	\textbf{Evidence (Multi-Metric Validation) :}
	\begin{itemize}
		\item \textbf{Brown Ocean dominated Long-lived tracks (2,3,5,6) :} Mean FSI = \textbf{0.42}; Mean ACF = \textbf{0.44}.
		\item \textbf{Runoff dominated Short-lived tracks (1,4) :} Mean FSI = \textbf{0.07}; Mean ACF = \textbf{-0.06}.
	\end{itemize}

	\vspace{0.1cm}
	\begin{center}
		\begin{tabular}{lcccc}
			\toprule
			\textbf{Track} & \textbf{TP-SM Lag} & \textbf{ACF-24h} & \textbf{Coupling}    & \textbf{Duration} \\
			\midrule
			1              & $-6h$              & $-0.061$         & Negative             & 75h               \\
			4              & $-12h$             & $0.032$          & Negative             & 65h               \\
			2, 3, 6        & $0$ to $+3h$       & $0.35-0.45$      & \textbf{Positive}    & 105-124h          \\
			5              & $+6h$              & $0.584$          & \textbf{Very Strong} & 163h              \\
			\bottomrule
		\end{tabular}
	\end{center}

	\vspace{0.1cm}
	\footnotesize
	\textbf{Interpretation :}
	\begin{itemize}
		\item \textbf{Positive Lag \& High ACF (Track 5) :} Soil moisture "leads" the system; high memory creates a self-sustaining survival strategy.
		\item \textbf{Zero Lag \& Moderate ACF (2,3,6) :} Efficient feedback where rainfall immediately recharges the moisture source.  \cite{osuri2020}
		\item \textbf{Negative Lag \& Low ACF (1,4) :} Rainfall does not translate to sustained soil moisture; energy source collapses post-landfall.
	\end{itemize}
\end{frame}

\subsection{Hypothesis 2}
\begin{frame}{Hypothesis 2 : Landfall Energy Budget Transition [1/2]}
	\vspace{-0.1cm}
	\begin{block}{The ocean-to-land transition creates an abrupt "Energy Gap" where MD survival depends on the land's ability to switch from a \texttt{LHF} dominant region to a Diurnal \texttt{SHF} amplified region}\end{block}

	\vspace{0.15cm}
	\footnotesize
	\textbf{Statement :} The ocean-to-land transition produces a significant reduction in surface evaporative cooling ($\Delta$THF $\approx +62.4$ W/m$^2$)\footnotemark[7]. Survival is governed by whether \textbf{Antecedent Soil Moisture} can buffer this large reduction in energy and sustain the core through enhanced sensible heating.
	% TODO:  Literature support : [Coupling between Land–Ocean–Atmosphere and Pronounced Changes in Atmospheric/Meteorological Parameters Associated with the Hudhud Cyclone of October 2014](https://pmc.ncbi.nlm.nih.gov/articles/PMC6313703/) and [Influence of Land Cover and Soil Moisture based Brown Ocean Effect on an Extreme Rainfall Event](https://www.nature.com/articles/s41598-019-53031-6)

	\vspace{0.2cm}
	\footnotesize
	$\Rightarrow$ \textbf{Landfall Signature :} \texttt{LHF}$\downarrow$ magnitude collapses while \texttt{SHF}$\uparrow$ becomes the diurnal engine. (Matches case study by \cite{andersen2014insights} but for Tropical Cyclone)

	\vspace{0.2cm}
	\footnotesize
	\textbf{The Transition Mechanism (energy shift happens via 2 level mechanism) :}
	\begin{enumerate}
		\item \textbf{Latent Collapse :} \texttt{THF} uptake drops sharply as the infinite moisture reservoir (Ocean) is cut-off. The system must immediately rely on stored terrestrial moisture (\texttt{SWVL1}).
		\item \textbf{Sensible Heating :} A transition from negative \texttt{SHF} over ocean ($-9.16$ W/m$^2$)\footnotemark[6] to positive \texttt{SHF} over land ($+4.50$ W/m$^2$)\footnotemark[6] provides the \textbf{diurnal instability} required for inland maintenance.
	\end{enumerate}

	% FIX:  Doesn't match in expampleblock
	% Track 1: Before -132.9 -> After -97.2 = +35.7 W/m$^2$
	% Track 4: Before -155.4 -> After -96.8 = +58.7 W/m$^2$
	% Track 5: Before -135.0 -> After -80.9 = +54.1 W/m$^2$
	\vspace{0.2cm}
	\begin{exampleblock}{\footnotesize The Spike or "Shock" Metric in the Heat flux change}
		\footnotesize Observed landfall flux transition : Magnitude change of (\textbf{$+68.8$ W/m$^2$})\footnotemark[6] (Before: $-153.6$ $\rightarrow$ After: $-84.8$ W/m$^2$).
	\end{exampleblock}

	\footnotetext[6]{Mean values are derived for Ocean-Origin tracks (1, 4 and 5), Land-Origin tracks have no values over ocean.}
	\footnotetext[7]{Upward energy supply from the surface has decreased by that amount.}
\end{frame}

\begin{frame}{Hypothesis 2 : Landfall Energy Budget Transition [2/2]}
	\footnotesize
	\textbf{Evidence (Survival vs. Dissipation Strategies) :}
	\begin{itemize}
		\item \textbf{Ocean-Origin Short-lived (Tracks 1, 4) :}
		      \begin{itemize}\footnotesize
			      \item \textbf{The Shock :} \texttt{THF} magnitude dropped by \textbf{\texttt{+41} to \texttt{+65 W/m$^2$}} after landfall. \cite{nair2011influence}
			      \item \textbf{The Failure :} Low soil capacity and Negative memory (\texttt{ACF < 0}) prevent the land from compensating for the loss of ocean evaporation. "Energy Gap" could not be filled; core moisture starved within \texttt{$\approx$70h}.
		      \end{itemize}
		\item \textbf{Ocean-Origin Long-lived (Track 5 - The Buffer) :}
		      \begin{itemize}\footnotesize
			      \item \textbf{The Shock :} Experienced the \textbf{largest} landfall shift (\texttt{$+70.5$ W/m$^2$}).
			      \item \textbf{Result :} High \textbf{Soil Moisture Memory (\texttt{ACF = 0.584})} acted as a thermodynamic buffer, sustaining the vortex until it reached the Indo-Gangetic Plain "Brown Ocean." This "Energy Gap" could be filled; thus core moisture sustained for \texttt{$\approx$163h}, despite a massive drop in immediate flux magnitude \texttt{$\Delta$ THF = +70.5 W/m$^2$} (Before landfall \texttt{$\approx$-155.3  W/m$^2$ $\rightarrow$ $\approx$-84.8 W/m$^2$} after landfall, \scriptsize 70.5 units closer to zero\footnotesize). \cite{nair2011influence}
		      \end{itemize}
		\item \textbf{Land-Origin Long-lived (Tracks 2, 3, 6) :}
		      \begin{itemize}\footnotesize
			      \item No discrete landfall shock; energy derived entirely from terrestrial recycling (\texttt{ACF-24h}). \cite{chang2009}
			      \item Represents the theoretical upper limit of coupling efficiency (\texttt{CE=15.35}) where land mimics ocean. \cite{hunt2017effect}
		      \end{itemize}
	\end{itemize}

	\vspace{0.1cm}
	\begin{center}
		\begin{tabular}{lccc}
			\toprule
			\textbf{Track Group} & \textbf{$\Delta$ Flux (W/m$^2$)} & \textbf{SHF Over Land} & \textbf{Outcome} \\
			\midrule
			Ocean (1, 4)         & \texttt{+53.3} (Mean)            & Weak/Broken            & Dissipation      \\
			Ocean + IGP (5)      & \texttt{+70.5}                   & Strong Diurnal         & Survival         \\
			Inland (2, 3, 6)     & $\approx 0$ (Mean)               & Sustained              & Persistence      \\
			\bottomrule
		\end{tabular}
	\end{center}
\end{frame}

\subsection{Hypothesis 3}
% Hyp 3 (ASM) is the Fuel in the tank. : The alsigns with thermodynamic principles => Wetter Antecedent Soil -> Greater moisture availability -> Enhanced evapotranspiration -> More latent heat release -> Increased atmospheric instability -> Stronger convection & precipitation feedback. ASM is a significant predictor of MD land-atmosphere coupling strength. The relationship is Statistically significant (p < 0.05) also correlation si very high r = 0.921.
% Track 3, 2024: Wettest soil (0.389) -> Strongest coupling (0.604)
% Track 1, 2024: Driest soil (0.198) -> Weakest coupling (0.027)
% r > 0.7 for both years
% - Works in dry years (2023) , r = 0.741, ASM = 0.169 (DRY), HF change at landfall +68.3 W/m2, TSI = 0.193
% - Works in wet years (2024) , r = 0.921, ASM = 0.316 (wet), HF change at landfall +26.3 W/m2, TSI = 0.193
% - TEST for year 2025)       , r = ?    , ASM = ?     (?  ), HF change at landfall +84.3 W/m2, TSI = ?
% Most depressions ALWAYS SHORT-LIVED (<96 hours)
% EVAPORATION is the key driver (latent heat dominates)
% H3 strength increases with wetter conditions
% The huge difference in soil moisture between 2023 and 2024 explains why H3 is stronger in 2024:
%   More soil moisture available
%   -> More evaporation feedback
%   -> Stronger ASM-FSI relationship
\begin{frame}{Hypothesis 3 : The Antecedent Moisture Paradox [1/2]}
	\begin{block}{Antecedent Soil Moisture (\texttt{ASM$_{24h}$}) enables initial infiltration, but it is \underline{\texttt{not}} a direct predictor of longevity. Duration is governed by Retention (Memory \texttt{ACF}) and Soil Type/Texture}\end{block}
	%
	% \vspace{0.15cm}
	% The antecedent soil moisture condition (24 hours pre-landfall) is a stronger predictor of monsoon depression behavior and energy budget characteristics than instantaneous precipitation during passage. Drier antecedent soils enhance infiltration capacity; wet antecedent conditions promote surface runoff and rapid downslope transport.

	\vspace{0.15cm}
	\footnotesize
	\textbf{Statement :} Pre-landfall soil moisture (\texttt{ASM$_{24h}$}) sets the state for landfall, but survival depends on whether the soil texture can translate that moisture into sustained fluxes. Without high memory ($ACF > 0.35$), the system collapses regardless of initial wetness.
	% \textbf{Statement :} Pre-landfall soil moisture (\texttt{$ASM_{24h}$}) sets the condition for the land-surface transition, but duration depends on whether the soil type (Clay over Indo-Gangetic) can maintain that moisture under the MD core. Without high memory (\texttt{ACF > 0.3}), even a saturated surface fails to sustain the Brown Ocean effect.

	\vspace{0.15cm}
	\footnotesize
	\textbf{The Infiltration vs. Runoff Logic :}
	\begin{itemize}
		\item \textbf{High ASM Failure:} In sandy/steep regions (Tracks 1, 4), high initial moisture leads to \textbf{Saturation Excess}. Rainfall cannot infiltrate further, triggering rapid surface runoff and downslope transport, breaking the feedback loop. \cite{kishtawal2013enhancement}
		\item \textbf{Low ASM Success:} In clay-rich plains (Track 5), even low initial ASM allows for high \textbf{Infiltration Capacity}. The soil "captures" the storm's precipitation, and high memory ($0.584$) keeps it available for evaporation.
	\end{itemize}

	\vspace{0.2cm}
	\footnotesize
	$\Rightarrow$ \textbf{Key Finding :} \texttt{ASM (> 0.03 m$^3$/m$^3$)} is a permissive condition but insufficient for longevity without positive \texttt{TP-SWVL1} coupling.

	\vspace{0.1cm}
	\begin{itemize}
		\item Ocean-origin tracks have pre-landfall \texttt{ASM} ranging from \texttt{0.028-0.071 m$^3$/m$^3$}, yet durations range from \texttt{75 - 165h}.
	\end{itemize}

	\vspace{0.1cm}
	\footnotesize
	\textbf{The "Paradox" Evidence :}
	\begin{itemize}
		\item \textbf{Track 4 (Shortest, 65h) :} Had the \textbf{highest} memory (\texttt{ASM = 0.071 m$^3$/m$^3$}) yet dissipated rapidly due to sandy soil over coast and moisture runoff by steep elevated terrain (Eastern Ghats). \cite{clapp1978empirical}
		\item \textbf{Track 5 (Longest, 163h) :} Had the \textbf{lowest} memory (\texttt{ASM = 0.028 m$^3$/m$^3$}). Survival was driven by the transition into the Indo-Gangetic Plain (with Clay soil with high moisture holding capacity) and avoiding/surviving the rapid runoff by high-elevation region (Eastern Ghats) by maintaining \texttt{CE > 10.5}.
	\end{itemize}

	\vspace{0.2cm}
	\footnotesize
	\textbf{Physical Logic :} High \texttt{ASM$_{24h}$} on sandy soil leads to \textbf{Saturation Excess or Runoff}, whereas moderate \texttt{ASM$_{24h}$} on clay soil leads to \textbf{Sustained Evaporation (\texttt{LHF})}.
	% TODO:  Literature support:
	% Emanuel & Sobel (2013): Land-atmosphere coupling critical for tropical cyclone intensity
	% Koster et al. (2004): Soil moisture influences precipitation patterns
	% Findell & Eltahir (2003): Feedback between land and atmosphere in tropical regions
\end{frame}

\begin{frame}{Hypothesis 3 : The Antecedent Moisture Paradox [2/2]}
	\footnotesize
	\textbf{Evidence (State vs. Process) :}
	\begin{itemize}
		% FIX: check needed
		% \item \textbf{Statistical Decoupling :} There is nearly zero correlation between $ASM_{24h}$ and total duration ($r \approx 0$).
		\item \textbf{Recharge Uniformity :} All tracks show a significant \texttt{SWVL1} increase (mositure availability) and have valid \texttt{ASM > 0.05 m$^3$/m$^3$} at landfall, but only Memory(\texttt{ACF > 0.3}) tracks convert this into multi-day persistence.
	\end{itemize}

	\vspace{0.1cm}
	\begin{center}
		\begin{tabular}{lcccc}
			\toprule
			\textbf{Track} & \textbf{ASM (24h)}   & \textbf{Soil Memory}       & \textbf{Soil Type + terrain}         & \textbf{Duration} \\
			\midrule
			1              & \texttt{0.037}       & \texttt{-0.06}             & Sandy + Steep                        & \texttt{75h}      \\
			4              & \texttt{0.071 (Max)} & \texttt{-0.15}             & Sandy + steep                        & \texttt{65h}      \\
			5              & \texttt{0.028 (Min)} & \underline{\texttt{+0.58}} & Mix(Sandy + Steep, then Clay + Flat) & \texttt{163h}     \\
			\bottomrule
		\end{tabular}
	\end{center}

	\vspace{0.1cm}
	\footnotesize
	\textbf{Interpretation :}
	\begin{itemize}
		\item \textbf{The Efficiency Gap :} \texttt{ASM$_{24h}$} provides the precondition but Soil Type act as reservoir for soil moisture (\texttt{ACF$_{24h}$}). Track 4 had a lot of \texttt{ASM} but a low \texttt{ACF}; Track 5 had low \texttt{ASM$_{24h}$} initially but a large \texttt{ACF$_{24h}$} that captured of precipitation (\texttt{TP}).
		\item \textbf{Thresholds :} An \texttt{ASM > 0.025 m$^3$/m$^3$}\footnotemark[7] is enough to initiate feedback, but the \texttt{ACF} and \textbf{Soil Type} determine how long it runs.
	\end{itemize}

	\vfill
	\footnotetext[7]{Threshold \texttt{ASM > 0.025 m$^3$/m$^3$} is the statistical low derived from Track 5, representing the minimum moisture state required to initiate terrestrial infiltration and coupling. While this sets the permissive boundary, survival is ultimately determined by the soil's retention capacity or memory (\texttt{ACF}) and texture rather than this initial value.}
\end{frame}

\subsection{Hypothesis 4}
% TODO:
%  - [ ] GAM (Generalized Additive Model) fit for LHF vs. SWVL1
%  - [ ] Second derivative plot showing curvature
%  - [ ] Identification of saturation onset threshold
% Hyp 4 (Coupling) is the Engine efficiency.
\begin{frame}{Hypothesis 4 : Nonlinear Moisture-Precipitation Dynamics [1/2]}
	\vspace{-0.1cm}
	\begin{block}{The \texttt{TP-SWVL1} relationship is fundamentally time-lagged. Instantaneous correlations ($r \approx 0.144$) hide the true feedback strength, which peaks at $+3$ to $+6$ hours post-rainfall}\end{block}

	\vspace{0.15cm}
	\footnotesize
	\textbf{Statement :} The peak soil (\texttt{SWVL1}) response lags precipitation (\texttt{TP}) by the infiltration timescale. Furthermore, the feedback sensitivity is governed by a \textbf{Saturation Plateau} in lag-correlation plot - drier soils show high sensitivity to rainfall pulses, whereas saturated soils reach a limit where additional rain only increases runoff.
	% \textbf{Statement :} The precipitation and soil moisture (TP--SWVL1) relationship is fundamentally nonlinear and time-lagged. Instantaneous correlations under-represent the feedback strength because the peak soil response lags precipitation by the infiltration timescale.
	% Statement : The relationship between precipitation and soil moisture is fundamentally nonlinear and time-lagged rather than instantaneous. A 6-12 hour lag exists between maximum precipitation intensity and peak soil moisture response, and the sensitivity of SWVL1 to TP changes is modulated by baseline antecedent saturation (drier soils show higher sensitivity; saturated soils show saturation-plateau behavior).

	\vspace{0.2cm}
	\footnotesize
	$\Rightarrow$ \textbf{Key Dynamic :} Low correlation instantaneous ($r=0.144$) hides peak coupling lags ($+3/+6h$); survival depends on avoiding the saturation cooling (saturation plateau in wet soils)

	\vspace{0.2cm}
	\footnotesize
	\textbf{Evidence :}
	\begin{itemize}
		%        \item Weak instantaneous correlation (r = 0.144) may reflect temporal mismatch
		\item \textbf{The Correlation Paradox :} Mean instantaneous correlation is low ($0.144$), but lag analysis reveals peak coupling at $+3$ to $+6$ hours for long-lived tracks.
		\item \textbf{The Saturation Plateau :} In high soil moisture Tracks(3 and 6), \texttt{SWVL1} plateaus at \texttt{0.42-0.45 m$^3$/m$^3$}. where additional TP does not increase SM but significantly increases \textbf{Runoff}, potentially cooling the surface and limiting FSI. \cite{clapp1978empirical}
		\item \textbf{Sensitivity Shift :} Drier soils for Track 2 and Track 5 (post-10th Sept, post-weakening after surviving elevated terrain of Eastern Ghats) shows the highest relative sensitivity, allowing small rainfall pulses to trigger large \texttt{LHF} and maintain feedback loop (\texttt{TP-SWVL1-LHF}). \cite{hunt2017effect}
	\end{itemize}
\end{frame}

\begin{frame}{Hypothesis 4 : Nonlinear Moisture-Precipitation Dynamics [2/2]}
	\footnotesize
	\textbf{Evidence (Temporal \& Physical Non-linearity) :}
	\begin{itemize}
		\item \textbf{Asymmetric Response :} Soil moisture rises rapidly during precipitation (Infiltration) but decays slowly (Evaporation(\texttt{LHF}) and Memory(\texttt{ACF})), creating the positive \texttt{ACF} seen in long-lived tracks.
		\item \textbf{Runoff vs. Feedback :} In Tracks 1 and 4, the absence of a lag-peak indicates that rainfall is lost to runoff before it can recharge the \texttt{SWVL1} layer. \cite{clapp1978empirical}
	\end{itemize}

	\vspace{0.1cm}
	% FIX: check from lag plots
	\begin{center}
		\begin{tabular}{lcccc}
			\toprule
			\textbf{System State}          & \textbf{Tracks} & \textbf{Correlation Type} & \textbf{Value (\texttt{r})} & \textbf{Physical Result}  \\
			\midrule
			Instantaneous                  & All (1--6)      & Linear Mismatch           & \texttt{0.144}              & Under-represents coupling \\
			\textbf{Lagged (\texttt{+6h})} & \textbf{2, 5}   & \textbf{Process-based}    & \texttt{0.45--0.62}         & \textbf{Feedback Peak}    \\
			Saturated                      & 3, 6            & Nonlinear Plateau         & \texttt{$\approx$0}         & Runoff / Cooling          \\
			Broken/Decoupled               & 1, 4            & Weak/Negative             & \texttt{< 0.1}              & Rapid Runoff              \\
			\bottomrule
		\end{tabular}
	\end{center}

	\vspace{0.3cm}
	\footnotesize
	\begin{itemize}
		\item Longevity is not just about having moisture in atmosphere; it is about the \textbf{efficiency} between a rainfall pulse and the soil response (\texttt{TP} $\rightarrow$ \texttt{SWVL1}).
	\end{itemize}
\end{frame}

\subsection{Hypothesis 5}
% FIX: update and refinement and furthur analysis required for this hypothesis
% Alternative Methodologies for H5: -> Implement lagged correlation analysis instead of Granger, Provide clearer picture of TP-SM temporal relationships (Already present)
% Use data from https://opentopography.org/
\begin{frame}{Hypothesis 5 : Two-Tier Control Mechanism}
	% TODO : Concept behind this hypothesis
	% - [x] Evidence A: Geographic Location. Your track data shows Tracks 3, 5, and 6 moved into the Indo-Gangetic Plains. Scientifically, this region is globally recognized as being dominated by Alluvial Clay/Silt. Tracks 1 and 4 are in Coastal/Ghats regions, which are known for Sandy/Loamy/Rocky soils.
	% - [x] Evidence B: The Autocorrelation Function (ACF). Your soilmoiture_analysis.py calculates Soil Memory.
	% - [x]     High Memory (ACF > 0.5): This is a physical "fingerprint" of Clay. Only Clay has the capillary strength to hold moisture that long.
	% - [x]     Low Memory (ACF < 0.2): This is a "fingerprint" of Sand.
	% - [x] Evidence C: The Saturation Plateau. In your coupling_metrics.py and correlation results, you saw SWVL1 cap out at 0.42–0.45. This specific value is the Porosity (Saturation Point) for typical Clay-Silt soils. Sandy soils usually saturate much lower (around 0.30–0.35).
	% - [ ] 1. H5 Soil: "Clay proxy ACF" -> Implicit. HWSD mandatory.
	% - [ ] 2. Orography r=-0.8?: Assumed -> compute SRTM.
	% - [ ] 3. CE Units: Inconsistent (0.5->15)-normalize properly.
	\vspace{-0.1cm}
	\begin{block}{MD longevity is governed by a hierarchical control: Soil Type (Tier 1) sets the capacity, while the Coupling Efficiency and ASM (Tier 2) determine the realization of that capacity}\end{block}
	% TODO: Mildly supporting Literature : [Observed increase in the peak rain rates of monsoon depressions](https://www.nature.com/articles/s41612-023-00436-w) and [Persistence in Physical Systems: An Application to Soil Moisture Memory](https://journals.aps.org/prl/abstract/10.1103/djz5-2lsn)

	\vspace{-0.5cm}
	\begin{alertblock}{\tiny Needs furthur analysis, maybe by using soil map data accurately finding Clapp-Hornberger ($\beta$)}\end{alertblock}

	\vspace{-0.3cm}
	% \vspace{0.15cm}
	\footnotesize
	\textbf{Statement :} The "Brown Ocean" effect is not a binary state but a two-tier mechanism. \textbf{Tier 1 (Geographical)}: Soil hydraulic properties (Clay vs. Sand) determine the theoretical lifespan. \textbf{Tier 2 (Dynamic)}: Pre-landfall moisture and feedback synchrony determine if that lifespan is achieved.

	\vspace{0.2cm}
	\tiny
	While I did not used a static soil texture map, the \texttt{ERA5-Land SWVL1} variable I analyzed is itself a product of the model's underlying soil parameters. By analyzing the Soil Memory (\texttt{ACF}) and the Saturation Plateau (\texttt{0.45 m$^3$/m$^3$}) in my results, I have identified the physical signatures of Clay-rich soils in the Indo-Gangetic Plains versus the high-drainage Sandy characteristics of the coastal tracks and steep elevation of Eastern Ghats.

	\vspace{0.05cm}
	\begin{block}{\footnotesize MD longevity is governed by a hierarchical control: Soil Type (Tier 1) sets the capacity \cite{seneviratne2010investigating}, while the Coupling Efficiency and \texttt{ASM} (Tier 2) determine the realization. \cite{hunt2017effect} \cite{kishtawal2013enhancement}.}\end{block}

	\vspace{0.05cm}
	\begin{enumerate}
		\item Clay (\texttt{$\beta$=0.1}): Mathematically, a low $\beta$ means the soil holds water tightly \cite{clapp1978empirical}.. For Tracks 3, 5, and 6 the soil moisture (\texttt{SWVL1}) stays high even after the rain stops (High Memory (\texttt{ACF})).
		\item Sand (\texttt{$\beta$=0.35}): A high $\beta$ means the soil drains water rapidly \cite{seneviratne2010investigating}. Tracks 1 and 4 where \texttt{SWVL1} drops almost immediately after a rain pulse.
	\end{enumerate}

	\vspace{0.1cm}
	\footnotesize
	\begin{itemize}
		\item \textbf{Tier 1: The Capacity (Necessary)}
		      \begin{itemize} \footnotesize
			      \item \textbf{Clay (IG Plains):} High retention ($\beta \approx 0.1$) $\rightarrow$ Capacity for \textbf{120--180h}.
			      \item \textbf{Sand (Coastal):} Rapid drainage ($\beta \approx 0.35$) $\rightarrow$ Capacity capped at \textbf{50--100h}.
		      \end{itemize}
		\item \textbf{Tier 2: The Efficiency (Sufficient)}
		      \begin{itemize} \footnotesize
			      \item \textbf{Coupling Synchrony:} Peak lag $0 \pm 6h$ + High ACF $\Rightarrow$ Maximum longevity.
			      \item \textbf{Failure Mode:} Negative lags + Broken loops $\Rightarrow$ Immediate decay regardless of soil type.
		      \end{itemize}
	\end{itemize}

	\vspace{-0.1cm}
	\footnotesize
	\begin{exampleblock}{\footnotesize Trajectory Validation}
		\textbf{Track 5:} Transitioned from Sandy Coast to IGP Clay $\rightarrow$ \textbf{163h} survival.\\
		\textbf{Track 1:} Remained in Coastal/Rocky terrain $\rightarrow$ \textbf{75h} dissipation.
	\end{exampleblock}
\end{frame}

\subsection{Hypothesis 6}
% Hyp 6 (Soil Type) is the Road Surface. -> if the road (soil) doesn't support the grip (moisture retention), the car (MD) won't go the distance.
% TODO :
% - HWSD/SoilGrids clay fraction at track coordinates
% - Clay% vs. Duration scatter plot
% - Clay% vs. ACF correlation
% FIX: Further analysis required for this hypothesis
\begin{frame}{Hypothesis 6 : Soil Type Modulation [1/2]}
	\begin{block}{Soil Type determines the efficiency of the "Land-act-as-Ocean" mechanism}\end{block}

	\footnotesize
	\textbf{Statement :} Soil type (Clay vs. Sand) determines the efficiency of the "Land-act-as-Ocean" mechanism by controlling moisture retention timescales \cite{seneviratne2010investigating}. Soil type narrates the efficiency of \texttt{TP} $\rightarrow$ \texttt{SWVL1} $\rightarrow$ \texttt{THF} energy transfer pathway \cite{hunt2017effect}. Clay soils enable longer feedback loops than sandy soils.

	\vspace{0.1cm}
	\begin{block}{\footnotesize Note :}
		Identification of soil type not via soil maps, but through \textbf{physical signatures} in ERA5-Land \texttt{SWVL1} variable.
	\end{block}

	\vspace{0.15cm}
	\footnotesize
	This Hypothesis uses Temporal Memory. By showing that Tracks 2, 3 5 and 6 have high \texttt{ACF >0.35}, proves that the soil itself is physically different (Clay-rich) from the Coastal tracks (Sand), where memory is near zero or negative.

	\begin{equation}
		\text{\texttt{TP}} \xrightarrow{\text{recharge}} \text{\texttt{SWVL1}} \xrightarrow{\text{flux}} \text{\texttt{LHF}} \xrightarrow{\text{evaporation}} \text{\texttt{TP}}
	\end{equation}
	Soil Texture (proxy (\texttt{ACF})) determines the efficiency of this entire chain. If the soil is sandy, the feedback loop (\cite{hunt2017effect}) breaks at the second link because the moisture infiltrates too deep too fast.

	\vspace{0.15cm}
	\footnotesize
	$\Rightarrow$ Clay retention $\rightarrow$ positive ACF/FSI.

	\vspace{0.15cm}
	\footnotesize
	\begin{center}
		\begin{tabular}{lllc}
			\toprule
			\textbf{Region} & \textbf{Soil Type} & \textbf{Tracks} & \textbf{Avg. Duration} \\
			\midrule
			Coastal/Delta   & Sandy              & 1, 4            & \textbf{70h}           \\
			Central Plains  & Clay/Loam          & 2, 3, 5, 6      & \textbf{126h}          \\
			\bottomrule
		\end{tabular}
	\end{center}

	\vspace{0.15cm}
	\footnotesize
	\textbf{Findings :}
	\begin{itemize} \tiny
		\item \textbf{Clay Dominance :} Indo-Gangetic tracks (Clay) exhibit positive SM-memory (\textbf{ACF $> 0.3$}) and high FSI (\textbf{$>0.4$}), facilitating feedback loops that double the MD lifespan.
		\item \textbf{Sandy Constraint :} Coastal tracks (Sand) suffer from \textbf{rapid infiltration} moisture vanishes from the top layer too quickly to sustain LHF, leading to negative autocorrelation.
	\end{itemize}

	% TODO: explicitly map soil type along each track using a soil dataset (e.g., FAO Soil Grids) to confirm clay vs. sandy regions.
	% \vspace{0.1cm}
	Clear regional pattern : clay soils over Indo-Gangetic plain support longer duration. (See Figure(\ref{fig:soilmapofindia}) Soil Type map of India)
\end{frame}

\begin{frame}{Hypothesis 6 : Soil Type Modulation [2/2]}
	\footnotesize
	\vspace{-0.1cm}
	\begin{columns}
		\begin{column}{0.45\textwidth}
			\textbf{Evidence A: Memory Signatures}
			\begin{itemize}
				\item \textbf{Clay Signature :} \texttt{ACF(24h) > 0.35}.
				\item Result : Track 5 (\texttt{0.584}) and Track 3 (\texttt{0.350}) show high persistence typical of the Indo-Gangetic Clay.
				\item \textbf{Sand Signature :} \texttt{ACF(24h) $\approx$ 0}.
				\item Result : Track 1 (\texttt{-0.061}) shows "Flash Infiltration" where moisture drains instantly.
			\end{itemize}
		\end{column}

		\begin{column}{0.55\textwidth}
			\textbf{Evidence B: Saturation Plateaus}
			\begin{itemize}
				\item \textbf{Clay Porosity :} \texttt{SWVL1 }saturates \texttt{0.42-0.45 m$^3$/m$^3$}. Matches tracks 2, 5, 6. Based on the $\beta$-parameter for tight soil retention, SWVL1 saturate at \texttt{0.45}. \cite{clapp1978empirical}.
				\item \textbf{Sandy Porosity :} \texttt{SWVL1} saturates lower (\texttt{< 0.35}). Matches coastal track 4.
			\end{itemize}
		\end{column}
	\end{columns}

	\vspace{0.2cm}
	\begin{equation*}
		\text{Clay Capacity (\texttt{160h+})} \gg \text{Sandy Capacity (\texttt{< 80h})}
	\end{equation*}

	\vspace{0.1cm}
	\footnotesize
	\textbf{ERA5 Correlation Findings :}
	\begin{itemize}
		\item \textbf{Feedback Efficiency :} Tracks 3, 5, and 6 transitioned into the Indo-Gangetic Plains. The high retention (Tier 1) allowed for sustained Latent Heat Flux even after rain pulses stopped.
		\item \textbf{The Failure Mode :} In Sandy/Rocky terrain (Track 1), the \texttt{TP} $\rightarrow$ \texttt{SWVL1} link is broken by rapid drainage, preventing the "Brown Ocean" from ever forming. \cite{wakefield2021}
	\end{itemize}
\end{frame}

\subsection{Hypothesis 7}
% TODO:
%   - [ ] SRTM 30m elevation data (lines 21-22 cite sources correctly)
%   - [ ] Elevation profiles along tracks
%   - [ ] Slope computation (°)
%   - [ ] HAND (Height Above Nearest Drainage)
%   - [ ] Correlation: Duration vs. Mean Elevation
\begin{frame}{Hypothesis 7 : Orographic Feedback Modulation}
	% Topographic sensitivity quantification for MD systems
	\vspace{-0.2cm}
	\begin{block}{The Topographic Constraint on the Brown Ocean Effect}\end{block}

	\footnotesize
	\textbf{Statement :} Elevation is the primary \textit{negative feedback modulator}. Even with optimal soil types (Hyp 6), terrain with elevation range \texttt{900-1600m} (Eastern Ghats) physically decouples the \texttt{TP} $\rightarrow$ \texttt{SWVL1} link via mechanical runoff dominance.

	\vspace{0.15cm}
	\footnotesize
	\textbf{The Three-Process Decay Mechanism:}
	\begin{enumerate} \footnotesize
		\item \textbf{Infiltration vs. Runoff :} Forced orographic ascent produces high-intensity rain pulses where \textbf{Precipitation is more than Infiltration Capacity}, forcing surface runoff.
		\item \textbf{Gravitational Drainage :} Steep slopes (avg $>20^\circ$ of Eastern Ghats) accelerate subsurface lateral flow, reducing the time for the top layer (\texttt{SWVL1}, 0-7 cm) to replenish mositurefrom precondition (\texttt{TP}).
		\item \textbf{LHF Starvation :} Rapid drainage leaves the surface effectively dry for evaporation to happen, thus the MD warm core lack latent heat and eventually dies off.
	\end{enumerate}

	\vspace{0.15cm}
	\footnotesize
	\begin{center}
		\begin{columns}[T]
			\begin{column}{0.48\textwidth}
				\textbf{Tier 1: High Elevation (range \texttt{900-1600 m})}
				\begin{itemize}
					\item \textbf{Tracks :} 1, 4 (Eastern Ghats)
					\item \textbf{Memory :} ACF $\approx 0$ (broken/oscillatory)
					\item \textbf{Longevity :} \textbf{Short (~70h)}
				\end{itemize}
			\end{column}
			\begin{column}{0.48\textwidth}
				\textbf{Tier 2: Low Plains (\texttt{$\approx$200 m})}
				\begin{itemize}
					\item \textbf{Tracks :} 2, 3, 5, 6 (Indo-Gangetic)
					\item \textbf{Memory :} ACF $> 0.35$ (sustained)
					\item \textbf{Longevity :} \textbf{Long (~130h+)}
				\end{itemize}
			\end{column}
		\end{columns}
	\end{center}

	\vspace{0.3cm}
	\centering
	\textbf{Max Elevation vs. Duration : Show negative correation} \\
\end{frame}

% \subsection{Hypothesis 8}
% FIX: Furthur analysis required for this hypothesis
% \begin{frame}{Hypothesis 8 : Topographic Resilience Hypothesis}
% \textbf{System Resilience During Topographic Barriers Depends on Upstream Soil Memory}
%
%      An MD crossing elevation > 900m will experience translational slowdown and temporary vorticity reduction. Survival depends on whether upstream soil moisture memory can sustain the TP-SWVL1 feedback loop during this critical period.
% Quantitative criterion:
%       If soil autocorrelation (24h) > 0.35 AND clay content > 40% AND TP-SWVL1 r > 0.3 -> Stall but survive If any metric falls below threshold -> Stall and die
%
%      Evidence:
%      Track 5: 0.584 autocorr, ~50% clay, r=0.44 -> Survives 165h despite 5->1 m/s slowdown
%      Track 1: -0.061 autocorr, ~30% clay, r=-0.15 -> Dies at 71h (no recovery)
%      Track 4: 0.032 autocorr, ~20% clay, r=0.20 -> Dies at 63h (no recovery)
%
% IN OBSERVATION :
% Why Vorticity Reduces Then Recovers :
% Vorticity = f(wind shear + thermal gradient + friction)
% During Ghats crossing :
%
%  - Wind shear reduced (blocked flow) -> vorticity decreases
%  - Friction increased (terrain drag) -> vorticity decreases
%  - Thermal energy sustained (TP-SWVL1 feedback active) -> vorticity doesn't decay completely
%
% Tracks 1,4 : Thermal energy lost -> vorticity continues decreasing -> eventual decay
% Track 5 : Thermal energy maintained -> vorticity stabilizes -> recovers as elevation decreases
%
% The Stall-Recover Mechanism:
%   + During stall: MD slows to 1 m/s, vorticity weakens temporarily
%   + If soil memory is strong: SM remains elevated -> TP-SWVL1 feedback sustains
%   + System stalls for 24-48 hours but doesn't decay
%   + Once elevation decreases: speed recovers, system re-energizes
% \end{frame}
